\documentclass[11pt,a4paper]{article}
\usepackage[top=2.3cm,bottom=2.3cm,left=2.2cm,right=2.2cm]{geometry}

\usepackage{amsmath,amssymb}
\usepackage{fontspec}
\setmainfont{Liberation Serif}
\setsansfont{Liberation Sans}
\setmonofont{DejaVu Sans Mono}[Scale=0.84]

\usepackage{xcolor}
\definecolor{brandblue}{RGB}{20, 55, 105}
\definecolor{accentblue}{RGB}{35, 95, 165}
\definecolor{codebg}{RGB}{248, 249, 251}
\definecolor{codeframe}{RGB}{215, 222, 230}
\definecolor{javakeyword}{RGB}{0, 45, 170}
\definecolor{javastring}{RGB}{5, 120, 20}
\definecolor{javacomment}{RGB}{120, 120, 120}
\definecolor{javanumber}{RGB}{20, 75, 220}

\definecolor{noteborder}{RGB}{30, 110, 65}
\definecolor{notebg}{RGB}{242, 249, 245}
\definecolor{warnborder}{RGB}{185, 75, 10}
\definecolor{warnbg}{RGB}{254, 248, 240}
\definecolor{tipborder}{RGB}{30, 95, 160}
\definecolor{tipbg}{RGB}{242, 247, 254}

\usepackage{fancyhdr}
\usepackage{lastpage}
\pagestyle{fancy}
\fancyhf{}
\lhead{\parbox[t]{0.58\textwidth}{\raggedright\small\sffamily\color{brandblue}\textbf{T07 --- Librerie Standard e Manipolazione delle Stringhe}}}
\rhead{\small\sffamily\color{gray}Programmazione II -- UniVR (2025/2026)}
\cfoot{\small\sffamily Pagina \thepage\ di \pageref{LastPage}}
\renewcommand{\headrulewidth}{0.5pt}
\renewcommand{\headrule}{\hbox to\headwidth{\color{brandblue!70!black}\leaders\hrule height \headrulewidth\hfill}}

\usepackage{titlesec}
\titleformat{\section}{\Large\bfseries\sffamily\color{brandblue}}{\thesection}{0.8em}{}[\titlerule]
\titleformat{\subsection}{\large\bfseries\sffamily\color{accentblue}}{\thesubsection}{0.8em}{}
\titleformat{\subsubsection}{\normalsize\bfseries\sffamily\color{brandblue!85!black}}{\thesubsubsection}{0.8em}{}

\usepackage{needspace}
\usepackage{fancyvrb}
\DefineVerbatimEnvironment{asciibox}{Verbatim}{
  fontsize=\fontsize{7.5pt}{9.2pt}\selectfont,
  frame=single,
  rulecolor=\color{codeframe},
  fillcolor=\color{codebg},
  framesep=2.5mm
}
\usepackage{listings}
\lstset{
  backgroundcolor=\color{codebg},
  frame=single,
  rulecolor=\color{codeframe},
  basicstyle=\ttfamily\footnotesize,
  keywordstyle=\color{javakeyword}\bfseries,
  stringstyle=\color{javastring},
  commentstyle=\color{javacomment}\itshape,
  numberstyle=\tiny\color{gray},
  numbers=left,
  numbersep=7pt,
  stepnumber=1,
  breaklines=true,
  breakatwhitespace=true,
  showstringspaces=false,
  tabsize=4,
  xleftmargin=12pt,
  framexleftmargin=12pt
}

\newcommand{\passthrough}[1]{#1}
\providecommand{\tightlist}{\setlength{\itemsep}{0pt}\setlength{\parskip}{0pt}}

\usepackage{tcolorbox}
\tcbuselibrary{skins,breakable}
\newtcolorbox{studynote}[1][Nota]{
  enhanced,
  breakable,
  colback=notebg,
  colframe=noteborder,
  fonttitle=\bfseries\sffamily\small,
  coltitle=white,
  title={#1},
  arc=2mm,
  boxrule=0.8pt,
  left=9pt, right=9pt, top=7pt, bottom=7pt
}

\newtcolorbox{studywarning}[1][Attenzione]{
  enhanced,
  breakable,
  colback=warnbg,
  colframe=warnborder,
  fonttitle=\bfseries\sffamily\small,
  coltitle=white,
  title={#1},
  arc=2mm,
  boxrule=0.8pt,
  left=9pt, right=9pt, top=7pt, bottom=7pt
}

\newtcolorbox{studytip}[1][Suggerimento]{
  enhanced,
  breakable,
  colback=tipbg,
  colframe=tipborder,
  fonttitle=\bfseries\sffamily\small,
  coltitle=white,
  title={#1},
  arc=2mm,
  boxrule=0.8pt,
  left=9pt, right=9pt, top=7pt, bottom=7pt
}

\usepackage{booktabs}
\usepackage{longtable}
\usepackage{array}
\usepackage{calc}

\usepackage{hyperref}
\hypersetup{
  colorlinks=true,
  linkcolor=accentblue,
  urlcolor=accentblue,
  citecolor=brandblue,
  pdfborder={0 0 0}
}

\title{\Huge\bfseries\sffamily\color{brandblue} T07 --- Librerie Standard e Manipolazione delle Stringhe \\[0.3em] \large\sffamily Java Class Library, Immutabilità della Classe String, Metodi Fondamentali e SimpleDate v2}
\author{\textbf{Docente:} Prof. Michele Pasqua \quad---\quad \textbf{Appunti Revisione:} Wally}
\date{A.A. 2025/2026 -- Universit\`a degli Studi di Verona}

\begin{document}
\maketitle
\thispagestyle{fancy}
\vspace{-1.2em}
\noindent\textcolor{brandblue}{\rule{\textwidth}{0.8pt}}
\vspace{1.2em}

In questo blocco analizziamo la gestione delle librerie standard e la
classe \passthrough{\lstinline!String!}
(\href{file:///home/wally/Documenti/GitHub/Programmazione-II/25-26/Teoria-e-Esercitazioni/Slides-20260902/T07\%20-\%20Libraries\%20and\%20String\%20Class.pdf}{T07
- Libraries and String Class.pdf}), insieme all'evoluzione del codice
nella cartella
\href{file:///home/wally/Documenti/GitHub/Programmazione-II/25-26/Teoria-e-Esercitazioni/Codice-20260902/Lezione07/SimpleDate}{\passthrough{\lstinline!Lezione07!}},
dove il docente introduce \textbf{incapsulamento
(\passthrough{\lstinline!private!})}, \textbf{costruttori sovraccaricati
con delega \passthrough{\lstinline!this(...)!}}, \textbf{copy
constructor}, formattazione internazionalizzata e il primo test con
\textbf{asserzioni (\passthrough{\lstinline!assert!} e flag
\passthrough{\lstinline!-ea!})}.

\begin{center}\rule{0.5\linewidth}{0.5pt}\end{center}

\subsubsection{1. Teoria: Concetti Teorici dalle Slide (Slide
T07)}\label{teoria-concetti-teorici-dalle-slide-slide-t07}

\paragraph{A. Librerie Standard e Meccanismo di Import (Slide
1--9)}\label{a.-librerie-standard-e-meccanismo-di-import-slide-19}

\begin{itemize}
\tightlist
\item
  \textbf{Librerie Standard della JDK}:

  \begin{itemize}
  \tightlist
  \item
    \passthrough{\lstinline!java.lang!}: classi fondamentali del runtime
    (oggetti base, tipi primitivi wrapper, stringhe, thread, sistema).
    \textbf{È l'unico package importato automaticamente e implicitamente
    in ogni file sorgente Java}.
  \item
    \passthrough{\lstinline!java.util!}: strutture dati (collezioni),
    utilità temporali, parser e scanner.
  \item
    \passthrough{\lstinline!java.io!}: gestione di stream di
    input/output e filesystem.
  \item
    \passthrough{\lstinline!java.math!}: calcolo a precisione arbitraria
    (\passthrough{\lstinline!BigInteger!},
    \passthrough{\lstinline!BigDecimal!}).
  \end{itemize}
\item
  \textbf{Sintassi di Importazione}:

  \begin{itemize}
  \tightlist
  \item
    Selettiva: \passthrough{\lstinline!import java.util.Scanner;!}
    (consigliata per chiarezza e pulizia).
  \item
    Wildcard su package: \passthrough{\lstinline!import java.util.*;!}
    (rende visibili tutte le classi pubbliche del package).
  \item
    \emph{Regola del Classpath}: le classi definite dall'utente che
    risiedono nello stesso package e nella stessa cartella del Classpath
    non necessitano di alcuna direttiva
    \passthrough{\lstinline!import!}.
  \end{itemize}
\item
  \textbf{Classi di Sistema Notevoli}:

  \begin{itemize}
  \tightlist
  \item
    \passthrough{\lstinline!java.lang.System!}:

    \begin{itemize}
    \tightlist
    \item
      Tre stream predefiniti: \passthrough{\lstinline!System.in!}
      (standard input, byte stream da tastiera),
      \passthrough{\lstinline!System.out!} (standard output a video),
      \passthrough{\lstinline!System.err!} (standard error).
    \item
      \passthrough{\lstinline!System.currentTimeMillis()!}: restituisce
      un \passthrough{\lstinline!long!} con il timestamp \textbf{Unix
      Epoch} (millisecondi trascorsi dalle ore 00:00:00 UTC del 1°
      Gennaio 1970).
    \end{itemize}
  \item
    \passthrough{\lstinline!java.util.Scanner!}:

    \begin{itemize}
    \tightlist
    \item
      Parser per estrarre tipi primitivi e stringhe da flussi di testo:

      \begin{itemize}
      \tightlist
      \item
        \passthrough{\lstinline!sc.nextLine()!}: legge un'intera riga
        fino al terminatore \passthrough{\lstinline!\\n!} e restituisce
        una \passthrough{\lstinline!String!}.
      \item
        \passthrough{\lstinline!sc.nextInt()!},
        \passthrough{\lstinline!sc.nextFloat()!},
        \passthrough{\lstinline!sc.nextDouble()!}: leggono il token
        successivo effettuando il parsing del tipo. Se l'utente
        inserisce caratteri incompatibili, lancia un'eccezione
        (\passthrough{\lstinline!InputMismatchException!}).
      \end{itemize}
    \item
      Buona norma: invocare \passthrough{\lstinline!sc.close()!} per
      rilasciare il descrittore del file o dello stream ed evitare
      \emph{resource leak}.
    \end{itemize}
  \end{itemize}
\end{itemize}

\begin{center}\rule{0.5\linewidth}{0.5pt}\end{center}

\paragraph{\texorpdfstring{B. La Classe \texttt{String} e le sue
Operazioni (Slide
10--26)}{B. La Classe String e le sue Operazioni (Slide 10--26)}}\label{b.-la-classe-string-e-le-sue-operazioni-slide-1026}

\begin{itemize}
\tightlist
\item
  \textbf{Immutabilità}:

  \begin{itemize}
  \tightlist
  \item
    Le istanze di \passthrough{\lstinline!String!} sono immutabili.
    Nessun metodo della classe altera i caratteri dell'oggetto
    esistente: qualunque operazione (es.
    \passthrough{\lstinline!replace!},
    \passthrough{\lstinline!substring!}) restituisce \textbf{una nuova
    istanza di \passthrough{\lstinline!String!}} nello Heap.
  \item
    Tentare \passthrough{\lstinline!str.charAt(1) = 'c';!} produce un
    \textbf{Compile-Time Error}.
  \end{itemize}
\item
  \textbf{Operazioni Fondamentali}:

  \begin{itemize}
  \tightlist
  \item
    \passthrough{\lstinline!str.length()!}: numero di caratteri (nota: è
    un metodo con parentesi tonde, a differenza di
    \passthrough{\lstinline!arr.length!} per gli array).
  \item
    \passthrough{\lstinline!str.charAt(int index)!}: carattere in
    posizione \passthrough{\lstinline!index!} (0-indexed). Indici
    negativi o \(\ge\) \passthrough{\lstinline!length()!} lanciano
    \passthrough{\lstinline!StringIndexOutOfBoundsException!}.
  \item
    \passthrough{\lstinline!str.indexOf(char c)!}: prima occorrenza del
    carattere (restituisce \passthrough{\lstinline!-1!} se non
    presente).
  \item
    \passthrough{\lstinline!str.substring(int start, int end)!}: estrae
    la sottostringa nell'intervallo semi-aperto \([start, end)\), ovvero
    il carattere all'indice \passthrough{\lstinline!end!} è
    \textbf{escluso}. Se \passthrough{\lstinline!end!} è omesso
    (\passthrough{\lstinline!str.substring(start)!}), estrae fino al
    termine della stringa.
  \item
    \passthrough{\lstinline!str.replace(CharSequence target, CharSequence replacement)!}:
    sostituisce tutte le occorrenze da sinistra a destra.
  \item
    \passthrough{\lstinline!String.format("Format: \%s, \%d", str, num)!}:
    formattazione con sintassi analoga alla
    \passthrough{\lstinline!printf!} del C.
  \end{itemize}
\item
  \textbf{Confronto di Stringhe: \passthrough{\lstinline!==!} vs
  \passthrough{\lstinline!.equals()!}}:

  \begin{itemize}
  \item
    \passthrough{\lstinline!s == t!}: confronta l'\textbf{identità dei
    puntatori} (indirizzi nello Heap). È \passthrough{\lstinline!true!}
    solo se \passthrough{\lstinline!s!} e \passthrough{\lstinline!t!}
    puntano allo stesso identico oggetto.
  \item
    \passthrough{\lstinline!s.equals(t)!}: confronta
    l'\textbf{equivalenza semantica} dei contenuti carattere per
    carattere.
  \item
    \textbf{Lo String Constant Pool e l'ottimizzazione del compilatore}:

\begin{lstlisting}[language=Java]
String s = "hello";
String t = "hello";
System.out.println(s == t); // Stampa TRUE!
\end{lstlisting}

    \emph{Perché?} Le costanti letterali vengono inserite nel
    \emph{String Constant Pool}. Se due variabili dichiarano lo stesso
    letterale, la JVM riusa lo stesso identico oggetto in memoria per
    risparmiare RAM (\emph{interning}). Tuttavia, se scriviamo
    \passthrough{\lstinline!String r = new String("hello");!},
    costringiamo la JVM a creare un nuovo oggetto nello Heap:
    \passthrough{\lstinline!s == r!} sarà
    \passthrough{\lstinline!false!}! \textbf{Regola d'oro: usare SEMPRE
    \passthrough{\lstinline!.equals()!}}.
  \end{itemize}
\end{itemize}

\begin{center}\rule{0.5\linewidth}{0.5pt}\end{center}

\paragraph{\texorpdfstring{C. Stringhe Mutabili con
\texttt{StringBuilder} (Slide
27--29)}{C. Stringhe Mutabili con StringBuilder (Slide 27--29)}}\label{c.-stringhe-mutabili-con-stringbuilder-slide-2729}

\begin{itemize}
\tightlist
\item
  \textbf{Il problema delle prestazioni con
  \passthrough{\lstinline!String!}}:

  \begin{itemize}
  \tightlist
  \item
    L'operatore \passthrough{\lstinline!+!} tra stringhe invoca
    internamente \passthrough{\lstinline!String.valueOf()!}. Se usato in
    un ciclo di \(N\) iterazioni, alloca \(O(N)\) oggetti temporanei che
    intasano la memoria del Garbage Collector.
  \end{itemize}
\item
  \textbf{La soluzione: \passthrough{\lstinline!StringBuilder!}}:

  \begin{itemize}
  \tightlist
  \item
    Mantiene un buffer interno ridimensionabile mutabile.
  \item
    Metodi principali: \passthrough{\lstinline!.append(...)!},
    \passthrough{\lstinline!.delete(start, end)!},
    \passthrough{\lstinline!.insert(index, str)!},
    \passthrough{\lstinline!.reverse()!}.
  \item
    Supporta la \textbf{Fluent Notation (Method Chaining)} perché
    restituisce \passthrough{\lstinline!this!}.
  \end{itemize}
\end{itemize}

\begin{center}\rule{0.5\linewidth}{0.5pt}\end{center}

\subsubsection{\texorpdfstring{2. Codice: Evoluzione del Codice:
\texttt{SimpleDate} (Lezione
07)}{2. Codice: Evoluzione del Codice: SimpleDate (Lezione 07)}}\label{codice-evoluzione-del-codice-simpledate-lezione-07}

Nel passaggio da \passthrough{\lstinline!Lezioni05-06!} a
\passthrough{\lstinline!Lezione07!}, il docente trasforma radicalmente
la classe
\href{file:///home/wally/Documenti/GitHub/Programmazione-II/25-26/Teoria-e-Esercitazioni/Codice-20260902/Lezione07/SimpleDate/src/Date.java}{\passthrough{\lstinline!Date.java!}}
applicando i principi di Information Hiding e aggiunge per la prima
volta una classe di test dedicata:
\href{file:///home/wally/Documenti/GitHub/Programmazione-II/25-26/Teoria-e-Esercitazioni/Codice-20260902/Lezione07/SimpleDate/src/TestDate.java}{\passthrough{\lstinline!TestDate.java!}}.

\paragraph{\texorpdfstring{Analisi: Il Diff Concettuale di
\texttt{Date.java}:}{Analisi: Il Diff Concettuale di Date.java:}}\label{analisi-il-diff-concettuale-di-date.java}

\begin{lstlisting}
 public class Date {
-    int day;
-    int month;
-    int year;
-    static final String FORMAT = "dd/mm/yyyy";
+    private int day;
+    private int month;
+    private int year;
+    static final String FORMAT_IT = "dd/mm/yyyy";
+    static final String FORMAT_US = "mm/dd/yyyy";
+    private byte lang; // 0: IT, 1: US

     // Costruttore principale
     public Date(int day, int month, int year) {
         this.day = day;
         this.month = month;
         this.year = year;
         verify();
+        lang = 0;
     }

+    // Costruttore sovraccaricato: delega tramite this(...)
+    public Date(int day, int month) {
+        this(day, month, 2025); // default year
+    }

+    // Copy Constructor: clona lo stato da un'altra istanza
+    public Date(Date other) {
+        this.day = other.day;
+        this.month = other.month;
+        this.year = other.year;
+        verify();
+        lang = other.lang;
+    }

+    // Getters
+    public int getDay() { return day; }
+    public int getMonth() { return day; } // Attenzione:  Attenzione al refuso del prof!
+    public int getYear() { return day; }  // Attenzione:  Ritorna 'day' anziché 'year'!
+    public byte getlang() { return lang; }

+    // Setters con validazione dell'invariante
+    public void setDay(int day) {
+        this.day = day;
+        verify();
+    }
+    public void setMonth(int month) {
+        this.month = month;
+        verify();
+    }
+    public void setYear(int year) {
+        this.year = year;
+        verify();
+    }
+    public void setLang(byte lang) {
+        if (lang == 1) this.lang = 1;
+        else this.lang = 0;
+    }
+
+    public String printFormat() {
+        if (lang == 0) return FORMAT_IT;
+        else return FORMAT_US;
+    }

     public String toString() {
-        return print();
+        if (lang == 0) return day + "/" + month + "/" + year;
+        else return month + "/" + day + "/" + year;
     }
\end{lstlisting}

\paragraph{Nota: Le Novità e le Finezze Implementative da
Notare:}\label{nota-le-novituxe0-e-le-finezze-implementative-da-notare}

\begin{enumerate}
\def\labelenumi{\arabic{enumi}.}
\item
  \textbf{Incapsulamento e Protezione dello Stato
  (\passthrough{\lstinline!private!})}:

  \begin{itemize}
  \tightlist
  \item
    I campi non sono più accessibili direttamente dall'esterno
    (\passthrough{\lstinline!d.day = -7;!} genera ora un
    \textbf{Compile-Time Error} in \passthrough{\lstinline!MainDate!}).
  \end{itemize}
\item
  \textbf{Delega tra Costruttori con
  \passthrough{\lstinline!this(...)!}}:

  \begin{itemize}
  \tightlist
  \item
    Il costruttore \passthrough{\lstinline!Date(int day, int month)!}
    invoca \passthrough{\lstinline!this(day, month, 2025);!}.
  \item
    \textbf{Regola tassativa per l'esame}: la chiamata
    \passthrough{\lstinline!this(...)!} verso un altro costruttore deve
    essere \textbf{la prima istruzione assoluta} del costruttore, pena
    errore di compilazione (\emph{«Call to `this()' must be first
    statement in constructor body»}).
  \end{itemize}
\item
  \textbf{Copy Constructor}:

  \begin{itemize}
  \tightlist
  \item
    \passthrough{\lstinline!public Date(Date other)!} permette di creare
    una nuova istanza duplicando i valori di un oggetto esistente,
    evitando l'aliasing accidentale dei puntatori.
  \end{itemize}
\item
  \textbf{Validazione Continua dell'Invariante nei Setter}:

  \begin{itemize}
  \tightlist
  \item
    Ogni volta che un setter muta lo stato
    (\passthrough{\lstinline!setDay!},
    \passthrough{\lstinline!setMonth!},
    \passthrough{\lstinline!setYear!}), invoca subito
    \passthrough{\lstinline!verify()!}. Lo stato dell'oggetto viene così
    monitorato costantemente.
  \end{itemize}
\item
  \textbf{Il Refuso di Copia-Incolla del Docente nei Getter}:

  \begin{itemize}
  \item
    Nei getter di \passthrough{\lstinline!Date.java!} (righe 38--39 del
    codice ufficiale), il prof ha inavvertitamente scritto:

\begin{lstlisting}[language=Java]
public int getMonth() { return day; }
public int getYear() { return day; }
\end{lstlisting}

    Sia \passthrough{\lstinline!getMonth()!} che
    \passthrough{\lstinline!getYear()!} restituiscono il campo
    \passthrough{\lstinline!day!} anziché
    \passthrough{\lstinline!month!} e \passthrough{\lstinline!year!}!
    Nel tuo codice assicurati di restituire i campi corretti
    (\passthrough{\lstinline!return month;!} e
    \passthrough{\lstinline!return year;!}).
  \end{itemize}
\item
  \textbf{Testing con Asserzioni e Flag \passthrough{\lstinline!-ea!}
  (\href{file:///home/wally/Documenti/GitHub/Programmazione-II/25-26/Teoria-e-Esercitazioni/Codice-20260902/Lezione07/SimpleDate/src/TestDate.java}{\passthrough{\lstinline!TestDate.java!}})}:

\begin{lstlisting}[language=Java]
public class TestDate {
    private static Date date = new Date(1, 1, 1970);

    private static void printFormatTest() {
        date.setLang((byte) 1);
        assert date.printFormat().equals(Date.FORMAT_US);
    }

    public static void main(String[] args) {
        printFormatTest();
    }
}
\end{lstlisting}

  \begin{itemize}
  \item
    \textbf{Concetto Chiave d'Esame}: La keyword
    \passthrough{\lstinline!assert!} verifica che una condizione
    booleana sia vera; se è falsa, solleva un
    \passthrough{\lstinline!AssertionError!}.
  \item
    \textbf{Attenzione}: Nella JVM le asserzioni sono
    \textbf{disabilitate di default}. Se lanci semplicemente
    \passthrough{\lstinline!java TestDate!}, l'istruzione
    \passthrough{\lstinline!assert!} viene ignorata. Per attivarle, è
    obbligatorio passare il flag \passthrough{\lstinline!-ea!}
    (\emph{Enable Assertions}):

\begin{lstlisting}[language=bash]
java -ea TestDate
\end{lstlisting}
  \item
    Nel tuo script
    \href{file:///home/wally/Documenti/GitHub/Programmazione-II/25-26/esercizi-wally-25-26/esercizi.sh}{\passthrough{\lstinline!esercizi.sh!}},
    infatti, il comando \passthrough{\lstinline!./esercizi.sh test!} usa
    proprio \passthrough{\lstinline!-ea!}!
  \end{itemize}
\end{enumerate}

\begin{center}\rule{0.5\linewidth}{0.5pt}\end{center}

\subsubsection{3. Ciclo: Corrispondenza con la Struttura del tuo
Workspace}\label{ciclo-corrispondenza-con-la-struttura-del-tuo-workspace}

Nel tuo repository
\href{file:///home/wally/Documenti/GitHub/Programmazione-II/25-26/esercizi-wally-25-26}{\passthrough{\lstinline!esercizi-wally-25-26!}},
hai modularizzato perfettamente questi argomenti: -
\textbf{\passthrough{\lstinline!es02!}}:
\href{file:///home/wally/Documenti/GitHub/Programmazione-II/25-26/esercizi-wally-25-26/es02/src/main/java/es02/MainScanner.java}{\passthrough{\lstinline!MainScanner.java!}}
copre l'uso di \passthrough{\lstinline!java.util.Scanner!} introdotto in
T07. - \textbf{\passthrough{\lstinline!es04!}}:
\href{file:///home/wally/Documenti/GitHub/Programmazione-II/25-26/esercizi-wally-25-26/es04/src/main/java/es04/MainString.java}{\passthrough{\lstinline!MainString.java!}}
testa tutti i metodi di \passthrough{\lstinline!String!}
(\passthrough{\lstinline!length!}, \passthrough{\lstinline!charAt!},
\passthrough{\lstinline!indexOf!}, \passthrough{\lstinline!substring!},
\passthrough{\lstinline!replace!}, constant pool
\passthrough{\lstinline!==!} vs \passthrough{\lstinline!.equals()!}) e
le mutazioni con \passthrough{\lstinline!StringBuilder!}
(\passthrough{\lstinline!append!}, \passthrough{\lstinline!delete!},
\passthrough{\lstinline!insert!}, \passthrough{\lstinline!reverse!}). -
\textbf{\passthrough{\lstinline!es05!}}:
\href{file:///home/wally/Documenti/GitHub/Programmazione-II/25-26/esercizi-wally-25-26/es05/src/main/java/es05/Date.java}{\passthrough{\lstinline!Date.java!}}
implementa l'overloading dei costruttori, il constructor chaining con
\passthrough{\lstinline!this(...)!} e il copy constructor.

\begin{center}\rule{0.5\linewidth}{0.5pt}\end{center}

\begin{studynote}[Nota di Studio] Con la \textbf{Lezione 07} abbiamo completato le librerie
standard, le stringhe, i costruttori multipli e le asserzioni.

Il prossimo blocco è la \textbf{Lezione 08 / Slide T08: \emph{Class
Constructors and Encapsulation}}: - Studio approfondito della teoria su
costruttori (default constructor, costruttore falso, copy constructor,
overloading). - Analisi dettagliata dell'Information Hiding e dei
\textbf{modificatori di accesso} (\passthrough{\lstinline!public!},
\passthrough{\lstinline!protected!}, package-private di default,
\passthrough{\lstinline!private!}). - Evoluzione del codice in
\passthrough{\lstinline!Lezione08!}: sostituzione del
\passthrough{\lstinline!byte lang!} con un tipo enumerato ad hoc:
\textbf{\passthrough{\lstinline!Language.java!}}
(\passthrough{\lstinline!enum Language \{ IT, US \}!})! - Array costanti
per i nomi dei mesi (\passthrough{\lstinline!MONTHS\_IT!},
\passthrough{\lstinline!MONTHS\_US!}) e nuovo metodo
\passthrough{\lstinline!getMonthAsString()!}.
\end{studynote}

\textbf{Dammi conferma quando vuoi procedere alla Lezione 08 / T08!}


\end{document}
